% =============================================================================
%  CHAPTER 6 -- RESULTS                                                [v3 OWNS]
% -----------------------------------------------------------------------------
%  v3.7 (2026-09-14): one section per environment, in the order of the study,
%  then a section that answers the two comparisons across all environments:
%    6.1 Obstacle avoidance       (DPCC's benchmark, 5 training seeds x 20 episodes)
%    6.2 Vision-conditioned alignment (seed 6, 10 paired contexts)
%    6.3 Quadrotor benchmark      (seed 6; corridor -> pillars -> s-curve)
%    6.4 Comparison across environments
%  Individual changelog: changelogs/v3.7_20260914_ch5-8_restructure.md
%
%  Labels referenced from v2's Chapter 4 and kept: sec:res:fewstep (6.1.2),
%  sec:res:constraints:degenerate (6.1.5). Others kept for v3's own references.
%  Vocabulary: Writing/Auxiliary/Naming/TRANSLATION_20260914_dev_jargon_to_scientific.md
% =============================================================================

\chapter{Results}\label{ch:results}

The results follow the order of the study: the obstacle-avoidance benchmark of \ac{DPCC} first, then
the two environments built for this thesis. \autoref{sec:res:summary} brings the three together and
answers, per environment, how the generative models and the projection methods compare. Each table
states its backbone and its evaluation protocol.

% =============================================================================
\section{Obstacle Avoidance}\label{sec:res:state}
% =============================================================================

All results in this section use five training seeds with twenty episodes each and are averaged per
geometry, then across the three geometries (\autoref{sec:setup:protocol:avoiding}).

\subsection{Generative Models against the Baseline}\label{sec:res:avoiding:models}

At a step budget of one, MeanFlow reaches a higher rate of success with constraint satisfaction than
the \ac{DPCC} baseline, in fewer control steps, at one thirty-first of its wall-clock time per control
step (\autoref{tab:state-headline}). Both use the same backbone with the same number of parameters.

\begin{table}[htpb]
  \caption[MeanFlow against the DPCC baseline]{MeanFlow against the DPCC baseline on the
  obstacle-avoidance task, per geometry and averaged over the three. Both use per-step projection on
  tightened constraints with four candidate plans, each with its own best selection rule; the baseline is
  at $\nfe=20$. Temporal U-Net, 4.0\,M parameters. Five training seeds $\times$ 20 episodes, except the
  baseline on both-hard, which has one training seed. S\&C: success with constraint satisfaction.}
  \label{tab:state-headline}
  \centering
  \small
  \begin{tabularx}{\linewidth}{@{}L rrrr@{}}
    \toprule
    Geometry & S\&C & Steps & ms/step & s/episode \\
    \midrule
    \multicolumn{5}{@{}l}{\emph{Diffusion (DPCC), $\nfe = 20$, cumulative projection cost}}\\
    top-left-hard   & 1.000 & 70.0 & 557.6 & 39.1 \\
    top-right-hard  & 0.950 & 77.6 & 517.1 & 40.2 \\
    both-hard       & 1.000 & 59.4 & 615.8 & 36.6 \\
    all three       & 0.983 & 69.0 & 563.5 & 38.9 \\
    \midrule
    \multicolumn{5}{@{}l}{\emph{MeanFlow, $\nfe = 1$, temporal consistency}}\\
    top-left-hard   & 1.000 & 58.9 & 18.4 & 1.08 \\
    top-right-hard  & 0.990 & 65.3 & 17.7 & 1.15 \\
    both-hard       & 0.990 & 58.8 & 18.2 & 1.07 \\
    all three       & 0.993 & 61.0 & 18.1 & 1.10 \\
    \midrule
    \multicolumn{5}{@{}l}{\emph{MeanFlow, $\nfe = 2$, temporal consistency}}\\
    top-left-hard   & 1.000 & 58.6 & 27.3 & 1.60 \\
    top-right-hard  & 0.980 & 64.6 & 27.3 & 1.76 \\
    both-hard       & 1.000 & 58.0 & 26.9 & 1.56 \\
    all three       & 0.993 & 60.4 & 27.1 & 1.64 \\
    \bottomrule
  \end{tabularx}
\end{table}

\dataref{temp/2508/batch\_avoiding\_combined\_20260825\_143212; DA\_20260827 \S10.1}

MeanFlow stays at or above $0.98$ on every geometry at both budgets, and the baseline reaches $0.950$ on
top-right-hard. Averaged over the three, the difference in success with constraint satisfaction, $0.993$
against $0.983$, lies within the across-seed standard error of about $0.07$; the difference in wall-clock
time is a factor of $31$ per control step and $35$ per episode. The baseline covers the both-hard
geometry with one training seed rather than five, which is why it has 220 episodes against 300.

Flow matching alone, at a step budget of two, reaches the success of the baseline on all three
geometries (\autoref{tab:state-fm}).

\begin{table}[htpb]
  \caption[Flow matching against the DPCC baseline]{Flow matching at $\nfe=2$ against the DPCC baseline
  at $\nfe=20$, per geometry, both with cumulative projection cost on tightened constraints. Temporal
  U-Net, 4.0\,M parameters; five training seeds $\times$ 20 episodes, except the baseline on both-hard
  (one training seed).}
  \label{tab:state-fm}
  \centering
  \small
  \begin{tabularx}{\linewidth}{@{}l L rrrr@{}}
    \toprule
    Geometry & Model & S\&C & Steps & s/step & s/episode \\
    \midrule
    top-left-hard  & Diffusion (DPCC) & 1.00 & 70.0 & 0.558 & 39.1 \\
                   & Flow matching    & 1.00 & 65.5 & 0.028 & 1.8 \\
    \midrule
    top-right-hard & Diffusion (DPCC) & 0.95 & 77.6 & 0.517 & 40.2 \\
                   & Flow matching    & 1.00 & 71.6 & 0.026 & 1.9 \\
    \midrule
    both-hard      & Diffusion (DPCC) & 1.00 & 59.4 & 0.616 & 36.6 \\
                   & Flow matching    & 1.00 & 68.0 & 0.028 & 1.9 \\
    \bottomrule
  \end{tabularx}
\end{table}

\dataref{DA\_20260819, $n=20$ comparison, \S1}

Flow matching needs fewer control steps than the baseline on the two single-halfspace geometries and
more on both-hard, where the baseline has one training seed. Its time per episode is 19 to 22 times
lower on all three. The success difference on top-right-hard, $1.00$ against $0.95$, is inside seed
noise, so what the table supports is that flow matching matches the success of the baseline at a
fraction of its time.

Run at the baseline's own budget of $\nfe=20$, the same flow-matching models reach the same success as
the baseline: paired by seed across three selection rules on both complete geometries, the difference over thirty
pairs is $-0.012$ with a standard deviation of $0.113$. At equal budget flow matching is $1.24$--$1.33$
times faster per control step, almost entirely in the projection; the generation cost per network
evaluation differs by nine per cent. Of the factor of twenty-one, about $1.4$ therefore comes from the
model at equal budget and about $15$ from running at a budget of two.
\dataref{DA\_20260819 \S1b, job 25006}

\autoref{fig:avoiding-tradeoff} places every configuration of the three models on the two cost axes.

\begin{figure}[htpb]
  \centering
  \includegraphics[width=\linewidth]{fig_avoiding_tradeoff}
  \caption[Control steps against time per control step]{Control steps against wall-clock time per
  control step, for each model at every step budget and both selection rules, on tightened constraints;
  temporal U-Net, 4.0\,M parameters, five training seeds $\times$ 20 episodes. The number next to a point
  is its step budget. The three transport configurations at the lower left reach the same number of
  control steps as the baseline at one twentieth to one thirtieth of its time per step.}
  \label{fig:avoiding-tradeoff}
\end{figure}

\subsection{Step Budget}\label{sec:res:fewstep}

The flow-based models keep their success with constraint satisfaction as the step budget falls to a
single network evaluation, while the diffusion baseline loses it below the budget it was trained with
(\autoref{fig:k-ladder}). At ten episodes per geometry the baseline reaches $0.667$ at $\nfe=1$ and
$1.000$ at $\nfe=10$ and $\nfe=20$; the flow-based models stay between $0.97$ and $1.00$ from $\nfe=1$ to
$\nfe=20$. For the diffusion model the number of steps is fixed when the noise schedule is discretised
at training time; for the flow-based models it is chosen by the ODE solver at inference.

\begin{figure}[htpb]
  \centering
  \includegraphics[width=\textwidth]{fig_avoiding_k_ladder}
  \caption[Success against step budget]{Success with constraint satisfaction against the step budget,
  averaged over the three geometries, temporal U-Net with 4.0\,M parameters. Solid: five training seeds
  $\times$ 20 episodes. Dashed, hollow markers: the baseline at five training seeds $\times$ 2 episodes.
  The baseline uses cumulative projection cost, the flow-based models temporal consistency.}
  \label{fig:k-ladder}
\end{figure}

\autoref{tab:state-lowk} gives, per geometry, the cheapest configuration of each model that reaches a
success rate with constraint satisfaction of $1.00$.

\begin{table}[htpb]
  \caption[Cheapest configuration at full success]{The cheapest configuration of each model that reaches
  S\&C $=1.00$, per geometry, chosen over all step budgets and all six selection rules by time per
  episode. Training seed 6, 20 episodes; temporal U-Net, 4.0\,M parameters; tightened constraints.}
  \label{tab:state-lowk}
  \centering
  \small
  \begin{tabularx}{\linewidth}{@{}L r L rr@{}}
    \toprule
    Model & $\nfe$ & Selection rule & Steps & s/episode \\
    \midrule
    \multicolumn{5}{@{}l}{\emph{top-left-hard}}\\
    Consistency training, $\alpha = 0.2$  & 1 & temporal consistency & 57.20 & 1.07 \\
    Consistency training, $\alpha = 0.05$ & 1 & temporal consistency & 58.20 & 1.09 \\
    MeanFlow                              & 1 & temporal consistency & 60.75 & 1.13 \\
    Flow matching                         & 1 & random & 64.40 & 1.30 \\
    Diffusion (DPCC)                      & 20 & cumulative projection cost & 61.00 & 35.66 \\
    \midrule
    \multicolumn{5}{@{}l}{\emph{top-right-hard}}\\
    Consistency training, $\alpha = 0.2$  & 2 & temporal consistency & 63.40 & 1.77 \\
    Consistency training, $\alpha = 0.05$ & 1 & cumulative projection cost & 73.15 & 1.23 \\
    MeanFlow                              & 5 & temporal consistency & 60.60 & 12.64 \\
    Flow matching                         & 1 & random & 70.55 & 1.21 \\
    Diffusion (DPCC)                      & 20 & cumulative projection cost & 61.35 & 29.48 \\
    \midrule
    \multicolumn{5}{@{}l}{\emph{both-hard}}\\
    Consistency training, $\alpha = 0.2$  & 1 & temporal consistency & 59.25 & 1.08 \\
    Consistency training, $\alpha = 0.05$ & 1 & random & 63.55 & 1.30 \\
    MeanFlow                              & 2 & temporal consistency & 58.90 & 1.58 \\
    Flow matching                         & 1 & temporal consistency & 64.55 & 1.27 \\
    Diffusion (DPCC)                      & 20 & cumulative projection cost & 59.35 & 36.55 \\
    \bottomrule
  \end{tabularx}
\end{table}

\dataref{temp/0309/batch\_avoiding\_combined\_20260903\_133730; Report\_20260903 \S1, \S3}

MeanFlow and consistency training reach full success in fewer control steps than flow matching on all
three geometries. In time per episode the geometries differ: on top-left-hard both are faster than flow
matching, while on top-right-hard MeanFlow reaches $1.00$ only at a budget of five, at $12.64$\,s per
episode against flow matching's $1.21$\,s, and on both-hard it needs a budget of two, at $1.58$\,s
against $1.27$\,s. Against the baseline in the same geometry, these configurations are 23 to 34 times
faster per episode on top-left-hard and both-hard, and between $2.3$ (MeanFlow) and 24 times faster on
top-right-hard.

\guard{MeanFlow and consistency training are not separated on this task:} on top-left-hard the step
difference between them has $p=0.231$ on one training seed, and averaged over all six selection rules
they are level, eight non-dominated configurations against nine.

\subsection{Plans before Projection}\label{sec:res:avoiding:raw}

Executing the plans without projection shows the generative models directly (\autoref{fig:raw-plans}).
At one network evaluation MeanFlow's successive plans form a narrow band around the executed path, each
a bounded curve towards the goal. The diffusion baseline at the same budget produces plans that jump
across the workspace from one control step to the next. Its executed path is nevertheless smooth,
because the tracking controller filters the plan.

\begin{figure}[htpb]
  \centering
  \begin{minipage}{0.49\linewidth}\centering
    \includegraphics[width=\linewidth]{fig_raw_plans_meanflow_K1}\\[0.3em]
    {\footnotesize MeanFlow, $\nfe=1$}
  \end{minipage}\hfill
  \begin{minipage}{0.49\linewidth}\centering
    \includegraphics[width=\linewidth]{fig_raw_plans_diffusion_K1}\\[0.3em]
    {\footnotesize Diffusion (DPCC), $\nfe=1$}
  \end{minipage}
  \caption[Plans without projection at one network evaluation]{Every plan generated along one episode,
  overlaid, executed without projection, \emph{both-hard}, training seed 6. Same backbone, same budget,
  and 0.0097 against 0.0094\,s per control step.}
  \label{fig:raw-plans}
\end{figure}

\dataref{Report\_20260819\_MF\_UNet \S7, figures 6a--6c}

Without projection, MeanFlow violates the constraints on $6.0$ control steps per episode and the
diffusion baseline on $28.0$, at $58.0$ and $56.5$ control steps. A second network evaluation brings
MeanFlow to $12.0$ with a narrower band of plans. On the same unprojected plans at $\nfe=1$ on
\emph{top-right-hard}, consistency training reaches the goal in 20 of 20 episodes, MeanFlow and flow
matching in 17 of 20, and the diffusion baseline at $\nfe=20$ in 12 of 20; with projection the difference
between consistency training and MeanFlow disappears.
\dataref{Report\_20260903\_AF\_UNet \S7; step counts over successful episodes}

\hole{Quantify plan smoothness --- jerk, path length or curvature over the saved plans --- to accompany
\autoref{fig:raw-plans}. The plans are on disk; no new runs are needed.}

\subsection{Projection Methods}\label{sec:res:avoiding:projection}

Endpoint projection \parencite{li2025hardflow} and per-step projection were run on the same MeanFlow
model, with the same solver and constraints, in one evaluation. Endpoint projection solves more
programs per control step and still takes about half the time: $0.0745$ against $0.1478$\,s at $\nfe=3$,
and $0.1304$ against $0.2268$\,s at $\nfe=5$ (\autoref{fig:projector-cost}). Its predicted endpoint is
close to the feasible set at every step, so the solver needs few iterations.

\begin{figure}[htpb]
  \centering
  \includegraphics[width=\textwidth]{fig_avoiding_projector_cost}
  \caption[Wall-clock time of the two projection methods]{Wall-clock time per control step of per-step
  and endpoint projection against the step budget; MeanFlow, temporal U-Net with 4.0\,M parameters,
  activation threshold 1.0, temporal consistency, tightened. The shaded budget has a single genuine step
  (\autoref{sec:res:constraints:degenerate}).}
  \label{fig:projector-cost}
\end{figure}

At $\nfe=3$, endpoint projection reaches a success rate with constraint satisfaction of $1.000$ with no
violations, in $58.88$ control steps at $0.0745$\,s per step. The \ac{DPCC} baseline reaches $1.000$ in
$70.13$ steps at $0.5534$\,s per step.
\dataref{temp/0609/I/batch\_avoiding\_combined\_20260906\_125724, job 25444; DA\_20260906 \S3}
\provisional{This comparison uses four training seeds (7--10) $\times$ three geometries $\times$ two
episodes, 24 episodes per configuration. The difference in success with per-step projection at the same
budget, $1.000$ against $0.917$, is two episodes and not significant. The evaluation at five training
seeds and twenty episodes is pending.}

\subsection{Genuine Steps}\label{sec:res:constraints:degenerate}

Endpoint projection needs sampling steps on which it is active and which are not the final step.
Their number is
\begin{equation*}
  n_{\text{gen}} \;=\; \max\!\big(\nfe - \lfloor (1-\eta)\,\nfe \rfloor,\, 1\big) - 1
\end{equation*}
(\autoref{sec:method:degenerate}). With $n_{\text{gen}}=0$ endpoint projection computes the same plan as
per-step projection, because at the final step the predicted endpoint is the current sample. The
evaluation code reports the count at run time: with the activation threshold at $1.0$ it ran with one
genuine step at $\nfe=2$, two at $\nfe=3$ and four at $\nfe=5$, and the number of solved programs grew as
$1:1.51:2.52$ against the predicted $1:1.5:2.5$. A budget of $\nfe=1$ has no genuine step at any
threshold. Configurations with $n_{\text{gen}}=0$ are reported separately and not averaged with others.
\dataref{job 25444 evaluation log, 2026-09-06}

\subsection{Candidate Plans and Selection Rules}\label{sec:res:ablations}

With a single candidate plan the three selection rules return the same plan: success and control steps
are identical in all fifteen paired seed--geometry combinations, for the diffusion baseline, MeanFlow,
flow matching and consistency training. With four candidates, generation time changes by at most six per
cent, because the four plans are one batched network evaluation, while projection time grows by a factor
of $3.2$ to $4.9$, because the programs are solved one after another on the CPU. Solving them in parallel
would bring MeanFlow at $\nfe=1$ from $18.1$ to about $11.7$\,ms per control step with all four
candidates kept; this is a prediction from two measured points.
\dataref{DA\_20260827 \S2, \S6, \S10.5}

Cumulative projection cost slows MeanFlow down on this task: $98.0$ control steps at $\nfe=2$ and $72.0$
at $\nfe=1$, against about $61$ with the other rules, at success rates of $0.927$ and $0.943$. The same
happens for consistency training and not for the diffusion baseline. Temporal consistency is the
selection rule for the flow-based models on this task.

% =============================================================================
\section{Vision-Conditioned Alignment}\label{sec:res:visual}
% =============================================================================

All results in this section use training seed 6 and the same ten contexts, and every difference is a
paired test (\autoref{sec:setup:protocol:aligning}). The box starts $0.4530$\,m from its target. No
configuration satisfies the full constraint set in every context, so the generative models are compared
on the final box-to-target distance of plans executed without projection, and constraint violations and
wall-clock time are compared separately.

\subsection{Generative Models}\label{sec:res:aligning:models}

MeanFlow brings the box closer to its target than the diffusion baseline in all ten contexts: median
final distance $0.0741$ against $0.4140$\,m, paired difference $-0.3744$\,m, $p=0.0020$, at $0.64$ of the
baseline's time per control step. It is also closer than flow matching, by $0.216$\,m, in nine contexts with none in favour of flow matching, $p=0.0039$; flow matching leaves the box unmoved in five contexts and MeanFlow in
one. Flow matching and the diffusion baseline end $0.0055$\,m apart, $p=1.0000$.
\dataref{Gen14 CLOSURE\_20260907 R.1, R.3}

Consistency training is level with MeanFlow at $\nfe=2$ and further from the target at $\nfe=20$. With
the final $\alpha=0.05$ it is worse in nine of ten contexts ($p=0.0215$); with $\alpha=0.2$ its mean
distance is $3.4$ times MeanFlow's, which the sign test does not reject.
\dataref{DA\_20260907 Gate 1; CLOSURE\_20260907 R.4}

\hole{Report the final orientation error of the box next to the final distance
(\texttt{context\_final\_box\_angle\_deg} against \texttt{context\_target\_angle\_deg}).}

\subsection{Step Budget}\label{sec:res:aligning:budget}

MeanFlow's median final distance falls from $0.2354$\,m at $\nfe=2$ to $0.0741$\,m at $\nfe=20$.
Consistency training does not improve with the budget: from $0.3738$ to $0.3918$\,m with $\alpha=0.2$,
and from $0.3978$ to $0.4488$\,m with $\alpha=0.05$. Its training target takes a finite difference from a
stop-gradient copy of the network, so it inherits part of that network's error, and that part does not
shrink as the step size does (\autoref{sec:method:alphaflow}).
\dataref{CLOSURE\_20260907 \S5.1}

\subsection{Projection Methods}\label{sec:res:aligning:projection}

At $\nfe=10$, endpoint projection matches the constraint satisfaction and the final distance of per-step
projection and saves $325$\,ms per control step, in all nine contexts that differ ($p=0.0039$, random
selection; $-320$\,ms and $p=0.0039$ with cumulative projection cost; $-211$\,ms and $p=0.0391$ with
temporal consistency). At $\nfe=20$ the two methods take the same time ($p$ between $0.18$ and $1.00$).
\dataref{CLOSURE\_20260907 \S8.1; H10 \S10.4}

MeanFlow at $\nfe=20$ with endpoint projection on tightened constraints keeps every context free of
violations, at $275.3$\,ms per control step, and moves the box in eight of ten contexts. The diffusion
baseline reaches the same freedom from violations at $2157.9$\,ms per control step; one of its variants
reaches it only by leaving the box unmoved in all ten contexts.
\dataref{H10 \S10.5}

\provisional{The diffusion baseline has not been evaluated on tightened constraints at $\nfe=20$, so the
constraint comparison against the baseline is incomplete. Repeated evaluations of configurations with
projection vary by about $0.4$\,m in final distance, so the minimum over contexts describes one episode.}

% =============================================================================
\section{Quadrotor Benchmark}\label{sec:res:uav}
% =============================================================================

The three scenes are reported in increasing difficulty. All results use training seed 6
(\autoref{sec:setup:protocol:uav}), and every generative model is a temporal U-Net with 3.96--3.97\,M
parameters.

\subsection{Corridor}\label{sec:res:uav:corridor}

Without projection every model flies into the slide on every flight: none of the twelve flights is
collision-free, with 42 to 56 violating control steps. With per-step projection at $\nfe=3$ and $\nfe=5$,
MeanFlow, consistency training and the diffusion baseline are collision-free on all twelve flights
($p=7\times10^{-7}$ against no projection).

The models differ in whether they complete the course. With projection, MeanFlow, consistency training
and flow matching cross the finish line on twelve of twelve flights and the diffusion baseline on none of
thirty-six: it stays collision-free but does not reach the line within the episode. Counting a flight
as a success when it crosses the finish line collision-free, MeanFlow and consistency training reach
$1.00$ at $\nfe=3$ and $\nfe=5$, flow matching at most $0.17$ because its flights are not collision-free,
and the diffusion baseline $0.00$ ($p=7\times10^{-7}$ for MeanFlow and consistency training against
either). Counting only flights that also reach the goal point, consistency training reaches $0.33$ and
MeanFlow $0.25$ at $\nfe=3$; after the slide only route L can reach its goal point, and these differences
are not significant. The flow-based models take $86$--$102$\,ms per control step at $\nfe=3$ and
$122$--$144$\,ms at $\nfe=5$; the projected diffusion baseline takes $626$--$651$\,ms.

Endpoint projection is $1.15$ to $2.1$ times faster than per-step projection on the same model and budget
--- $47$ and $76$\,ms against $99$\,ms for MeanFlow at $\nfe=3$ --- but collision-free on at most $0.58$ of
the flights, against $1.00$ ($5$ of $12$ against $12$ of $12$, $p=0.005$), with $0.6$ to $3.2$ violating
steps per flight.
\dataref{Gen15/U16/DA\_20260914\_corridor\_v2\_paper\_full.md \S0--\S7; temp/1409/batch\_uav\_20260914\_091148}
\guard{Twelve flights per configuration, four per route, one training seed. The diffusion model is trained
with action loss weight 1.}

\subsection{Pillars}\label{sec:res:uav:pillars}

Averaged over seventeen matched configurations at $\nfe=5$, the success rate with constraint satisfaction
is $0.635$ for MeanFlow, $0.359$ for flow matching and $0.235$ for consistency training. With the full
constraint set and projection, the best configurations reach $0.90$, $0.50$ and $0.70$; seven of
MeanFlow's seventeen configurations reach $0.8$ or more, two of flow matching's, and none of consistency
training's.
\dataref{Gen15 CLOSURE\_20260910 \S2; DA\_20260912 pillars reference \S5.1}

Endpoint projection takes $8.2$ to $21.3$ times less projection time than per-step projection with the
full geometry ($5.6$ to $12.4$ times less in total), at comparable success. None of its 210 flights left
the constraint set; 43 of the 270 flights with per-step projection did, and seven of the eight
configurations with such flights used random selection.
\dataref{Gen15 T3 \S3--\S6}

The diffusion baseline, trained at $\nfe=20$ with action loss weight 1, crosses the finish line on all
fifty flights, never comes within $0.30$\,m of the goal point, and violates the pillar clearance on 101 to
156 control steps per flight. Its success rate with constraint satisfaction is $0.00$ at the goal point
and at most $0.30$ at the finish line, at $4425$\,ms per control step. At $\nfe=5$ the flow-based models reach up to $0.70$--$1.00$ and $0.90$--$1.00$ on the same two criteria, counting configurations without projection, and up to $0.50$--$0.90$ at the goal point with full projection.
\dataref{DA\_20260912 pillars diffusion reference, rev.~2}
\provisional{Ten flights per configuration on one training seed. The diffusion budget is a training-time
choice and the flow-based budgets are inference-time choices, so the baseline appears here as a
reference. The multi-seed evaluation is pending.}

\subsection{S-Curve}\label{sec:res:uav:scurve}

The s-curve is the most difficult scene. Across forty configurations the success rate with constraint
satisfaction is at most $0.100$, and raising the step budget beyond five does not increase it. The
tracking controller does not follow the plans through the two sharp turns, and the order of the models
on this scene changes with the controller used, so no ordering of the models is drawn from it.
\dataref{Gen15 CLOSURE\_20260910 \S3, \S6; DA\_20260909 T5 controller comparison}

% =============================================================================
\section{Comparison across Environments}\label{sec:res:summary}
% =============================================================================

\subsection{Generative Models}\label{sec:res:summary:models}

\autoref{tab:summary-models} collects, per environment, the four comparisons between the generative
models. MeanFlow is ahead of flow matching and of the diffusion baseline in every environment.
Consistency training is level with MeanFlow on obstacle avoidance and on the corridor, and behind it on
alignment at $\nfe=20$ and on the pillars. Flow matching reaches the success of the baseline at a
fraction of its time on obstacle avoidance, ends level with it on alignment, and completes the corridor
where the projected baseline does not.

\begin{table}[htpb]
  \caption[Generative models across environments]{The generative models compared per environment. $>$:
  ahead with the stated evidence; $\approx$: not separated; $<$: behind. Temporal U-Net, 3.96--4.0\,M
  parameters, throughout.}
  \label{tab:summary-models}
  \centering
  \small
  \begin{tabularx}{\linewidth}{@{}p{2.1cm} L L L L@{}}
    \toprule
    Environment & Consistency training vs.\ MeanFlow & MeanFlow vs.\ flow matching & Flow matching vs.\ diffusion & MeanFlow vs.\ diffusion \\
    \midrule
    Obstacle avoidance (5 seeds $\times$ 20) & $\approx$ ($p=0.231$) & $>$ fewer steps on all three geometries; time per episode depends on the geometry (seed 6) & $\approx$ success, $21\times$ faster & $>$ S\&C, fewer steps, $31\times$ faster \\
    Alignment (seed 6, 10 contexts) & $\approx$ at $\nfe=2$; $<$ at $\nfe=20$ ($p=0.0215$) & $>$ $0.216$\,m, $p=0.0039$ & $\approx$ $p=1.0000$ & $>$ $0.374$\,m, $p=0.0020$ \\
    Corridor (seed 6, 12 flights) & $\approx$ & $>$ collision-free $12/12$ vs.\ $0/12$ & completes $12/12$ vs.\ $0/36$ & $>$ $1.00$ vs.\ $0.00$, $p=7\times10^{-7}$ \\
    Pillars (seed 6, 10 flights) & $<$ $0.235$ vs.\ $0.635$ & $>$ $0.635$ vs.\ $0.359$ & best configurations above the baseline & best configurations above the baseline \\
    \bottomrule
  \end{tabularx}
\end{table}

The order MeanFlow $>$ flow matching $>$ diffusion therefore holds in control steps on obstacle avoidance, in final distance on alignment, and in collision-free flights on the corridor. Consistency training is not ahead of MeanFlow in any
environment. Where the task separates them, its finite-difference target leaves part of the network's
own error in the field regardless of the step budget (\autoref{sec:method:alphaflow}).

\subsection{Projection Methods}\label{sec:res:constraints}

\autoref{tab:summary-projection} compares endpoint projection with the per-step projection of \ac{DPCC}
in each environment. Endpoint projection takes less time in all four. Its constraint satisfaction is
level with per-step projection on obstacle avoidance and on alignment, higher on the pillars, and lower on
the corridor.

\begin{table}[htpb]
  \caption[Projection methods across environments]{Endpoint projection against per-step projection, same
  generative model and budget. Temporal U-Net, 3.96--4.0\,M parameters.}
  \label{tab:summary-projection}
  \centering
  \small
  \begin{tabularx}{\linewidth}{@{}p{3.0cm} L L@{}}
    \toprule
    Environment & Time & Constraint satisfaction \\
    \midrule
    Obstacle avoidance, $\nfe=3$ (24 episodes) & $0.0745$ vs.\ $0.1478$\,s per step & S\&C $1.000$ vs.\ $0.917$, not significant \\
    Alignment, $\nfe=10$ (10 contexts) & $325$\,ms per step less, $p=0.0039$ & level, same final distance \\
    Corridor, $\nfe=3$ (12 flights) & $1.15$--$2.1\times$ less & collision-free $\le0.58$ vs.\ $1.00$, $p=0.005$ \\
    Pillars, $\nfe=5$ (10 flights) & $8.2$--$21.3\times$ less projection time & $0/210$ vs.\ $43/270$ flights violating \\
    \bottomrule
  \end{tabularx}
\end{table}

All endpoint-projection results above have at least two genuine steps
(\autoref{sec:res:constraints:degenerate}). \provisional{Every entry except obstacle avoidance rests on
one training seed.}
